The efficacy of the Deep Reinforcement Learning (DRL) strategy relied on its ability to exploit two underlying behavioral mechanisms uncovered during the simulation: the stabilizing force of ``Cultural Inertia'' and the opposing friction of ``Density-Dependent Resistance.'' Unlike the static policy regimes, the adaptive agent successfully identified these non-linear dynamics, leveraging social momentum to sustain compliance in rural areas while concentrating capital to overcome the structural inertia of urban centers \citep{Jimenez2025}.

\subsection{Cultural Inertia and the ``Graduation'' Effect}

A critical finding of this study was the pivotal role of Social Norms in sustaining compliance behavior even after the withdrawal of direct financial intervention. This phenomenon, observed in barangays such as Mati and Binuni, manifested as \textbf{Cultural Inertia}, where a community maintained high segregation rates despite the DRL agent cutting the local budget to near-zero ``Maintenance Mode'' levels. 

For instance, the simulation data revealed that Barangay Mati maintained 100\% compliance from Quarter 4 through Quarter 12, even as the agent reallocated the majority of its funding to other, more problematic areas. This persistence suggested that once a specific compliance threshold was crossed, the behavior became self-reinforcing, decoupling the action from the immediate presence of monetary rewards or enforcement risks \citep{Corcoran2020}.

This mechanism aligned with the theoretical framework established by \citet{Andre2021}, who argued that social norms act as a ``ratchet mechanism'' for pro-environmental behavior—hard to build, but resistant to decay once established. Furthermore, this validates the \textit{Threshold Hypothesis} in Philippine governance, where \citet{Nishimura2022} observed that collective action becomes sustainable only after a critical mass of the population internalizes the mandate, effectively allowing the LGU to ``graduate'' compliant barangays and redirect scarce resources elsewhere.

\subsection{Density-Dependent Resistance and the Urban Trap}

Conversely, the struggle in Barangay Poblacion revealed that resistance was driven by ``Density-Dependent'' factors. Poblacion's large population of 1,534 households meant that even when the DRL agent allocated 90\% of the total municipal budget to this single barangay, the resulting enforcement intensity only reached an initial effective value of approximately 0.38. 

This dilution effect confirmed the \textbf{Urban Resource Trap} described by \citet{Villanueva2021}, where high population density dissolved the impact of standard monitoring mechanisms, allowing non-compliance to persist as the status quo despite significant absolute spending. This phenomenon is consistent with global observations by \citet{Kaza2021}, noting that waste generation rates and enforcement difficulty scale non-linearly with urbanization.

The simulation implied that for such high-density urban centers, standard budget allocation strategies were insufficient to break the initial resistance threshold. The data suggested the necessity of a specialized ``Shock and Awe'' policy, where funds were temporarily diverted entirely from Information, Education, and Communication (IEC) and Incentives to focus purely on Enforcement. This aligned with the observations of \citet{Badua2022} regarding the operational burdens of Local Government Units, highlighting that overcoming the inertia of large, anonymous urban populations required a threshold of enforcement visibility that equal-distribution strategies could never mathematically achieve. The DRL agent's eventual success in Poblacion proved that overcoming this density barrier was possible, but only through the sequential concentration of resources to artificially induce the behavioral tipping point \citep{Centola2018}.